% TOP_PROBLEM: {"id": 1101303, "title": "Questions in Geometric Group Theory — Q 13.3", "category_id": 17, "category": {"id": 17, "name": "group_theory", "display_name": "Group Theory", "description": "Problems about groups, group actions, representations, and related algebraic structures.", "slug": "group-theory", "order_index": 17, "created_at": "2026-07-31T15:26:25.671Z"}, "status": "open", "problem_number": "AMR-010-1303", "set_id": 13}
% FIRST_POSTED: 2026-10-02
% ENTRY_KIND: statement_only
% UnsolvedMath ID 1101303; source catalogue code AMR-010-1303
% !TeX program = lualatex
% Definition and sources only; this file is not an LLM solution attempt.
\documentclass[11pt,a4paper]{article}
\usepackage[margin=25mm]{geometry}
\usepackage{fontspec}
\setmainfont{lmroman10-regular.otf}[BoldFont=lmroman10-bold.otf,
 ItalicFont=lmroman10-italic.otf,BoldItalicFont=lmroman10-bolditalic.otf]
\usepackage{amsmath,amssymb,mathtools}
\usepackage{unicode-math}
\setmathfont{latinmodern-math.otf}
\AtBeginDocument{\let\setminus\smallsetminus}
\usepackage{xurl,hyperref}
\hypersetup{colorlinks=true,urlcolor=blue}
\setcounter{secnumdepth}{0}
\setlength{\parindent}{0pt}
\setlength{\parskip}{5pt}
\setlength{\emergencystretch}{3em}
\begin{document}
\section*{Questions in Geometric Group Theory — Q 13.3}\label{1101303}
{\small UnsolvedMath ID 1101303 · AMR-010-1303 · Group Theory · AMR Open Problem Lists}
\subsection{Definitions and mathematical statement}
A finite graph is planar if it admits an embedding in the two-sphere in which edges meet only at common endpoints. For connected finite graphs, a covering map $p:H\to G$ is a graph map such that for each vertex $v\in V(H)$ the induced map on incident edges is a bijection onto the edges incident to $p(v)$. Its degree is the common finite cardinality of the vertex fibers. A double cover has degree two; no regularity assumption is imposed on a general finite cover.

Is the following implication true for every connected finite graph $G$?
\[
 \text{$G$ has a finite planar cover}
 \quad\Longrightarrow\quad
 \text{$G$ is planar or has a planar double cover.}
\]
Equivalently, the least degree of a planar finite cover, if any exists, is one or two. This is the source's $1$--$2$--$\infty$ formulation of the planar-cover problem.

The covering graph is not an arbitrary graph mapping onto $G$: the local bijection condition must hold at every vertex. In particular, merely drawing a subdivision or a branched cover does not establish a planar cover. One can formulate the question for simple graphs; loops and parallel edges require the same incidence-based covering convention. For a disconnected base graph, the issue is tested component by component, so the connected formulation contains the substantive problem.
\subsection{Short English statement}
If a finite graph has any finite planar cover, must it already have one of degree at most two?
\subsection{Sources}
\begin{itemize}
\item[S1] ULAM.AI and the UnsolvedMath Contributors, supplied problems.json, record 1101303 (AMR-010-1303), AMR Open Problem Lists. Catalogue identity and source transcription; CC BY 4.0.
\url{https://huggingface.co/datasets/ulamai/UnsolvedMath}.
\item[S2] Bestvina - Questions in Geometric Group Theory (pdf) (2004), Question 13.3 (PDF page 24). Original source indexed by AMR-010-1303.
\url{https://www.math.utah.edu/~bestvina/eprints/questions-updated.pdf}.
\end{itemize}
\end{document}
